\section{Introduction}
\label{sec:intro}
Metaheuristics that learn while they search are a prominent line of work in evolutionary computation and operations research. A typical \emph{learnheuristic}~\cite{CalvetEtal2017} embeds a reinforcement-learning agent, most often tabular Q-Learning~\cite{WatkinsDayan1992,SuttonBarto2018}, as an adaptive selector of operators, strategies or parameters inside a search framework such as Iterated Local Search (ILS)~\cite{LourencaoEtal2019} or Iterated Greedy (IG)~\cite{RuizStutzle2007}. The usual evidence is a comparison with the search without the agent, or with published baselines, on benchmark instances~\cite{KarimiMamaghanEtal2023,MissaouiOSullivan2026}; when the learned version wins, the result is read as evidence that the \emph{learned policy} is what helps.

That reading is fragile for three reasons. First, the \emph{search structure} that hosts the agent (which operators it can choose from, which acceptance rule it uses, whether it contains a local search) may produce most of the gain under any selection rule, so a comparison against a different search framework confounds the policy with the structure. Second, the \emph{budget}: an agent changes the CPU time per iteration, and a local search multiplies it, so equal iteration counts are not equal budgets; a policy that merely chooses cheaper actions buys extra iterations and looks better at equal time without having decided anything better. Third, the \emph{tuning}: a parameter validated offline on separate instances is a competitor that most studies omit, although a practitioner would use it before reaching for a learner. Adaptive operator selection (AOS) is a well-studied topic~\cite{FialhoEtal2010,PisingerRopke2007}, but we found it rarely answered with these controls: in an audit of $23$ recent open-access papers that embed a learned selector in a metaheuristic, $10$ include a random-selection control, $6$ compare at a common time limit, $3$ include a static configuration tuned offline, and $22$ report a benefit of the learned component ($5$ of them only partly). The audit is small and has a single coder; its coding sheet is in the repository.

\paragraph{Research questions} We ask three questions.
\textbf{RQ1:} what explains the advantage of a learning-guided IG over ILS, and how much of a local search's gain is only budget?
\textbf{RQ2:} do online learners beat random choice over the same actions, at equal iterations and at equal CPU time; does a static configuration win anyway; and could learning have shown up if it existed?
\textbf{RQ3:} where, if anywhere, does an online selector pay, can that be foretold, and what exploits the headroom instead?

\paragraph{Contributions}
(i) A \emph{budget-fair attribution protocol} (Section~\ref{sec:protocol}) built from a random control at a matched budget, a split of a learner's gain at equal CPU time into a decision effect and a cost effect, a static competitor validated offline on separate instances, a cross-fitted oracle headroom with a formal treatment of its bias (Section~\ref{sec:estimator} and Appendix~\ref{app:estimator}), analytic and planted-gain power checks, positive controls, and predictions recorded with code hashes before every battery (Section~\ref{sec:stats}); the protocol is released as a small library.
(ii) Evidence from $\nbatteries$ pre-specified batteries on four hosts (identical parallel machines, unrelated machines, a permutation flowshop, a genetic algorithm), on synthetic classes, on $13{,}525$ instances of a published benchmark of $P||\Cmax$, on three published families of $R||\Cmax$ instances and on the standard $R||\Cmax$ benchmark of Fanjul-Peyro and Ruiz with best-known solutions, that the advantage of the learning-guided IG comes from its search structure, that the share of a local search's gain that is only budget depends on the efficiency of the implementation, and that the standard Q-Learning wins at equal CPU time by choosing cheaper actions, not by a learned policy.
(iii) A \emph{headroom analysis}: the usual one-seed oracle headroom is mostly run-to-run noise; with the noise removed, the headroom is a necessary condition for Q-Learning to pay, though not a sufficient one; a policy learned \emph{offline} (a table by class, or a neural network over the state of the search) exploits it, mostly through class-level knowledge, and a validated static action or mixture was never beaten beyond the equivalence margin, including by irace-tuned learners, two published AOS designs, a neural policy, two state-dependent look-ahead policies that are given hundreds of times the evaluations of a learner, and a horizon ten times longer.
(iv) A candid account of a reversal: on published families of unrelated-machine instances with correlated machines the standard Q-Learning seemed to recover half to two thirds of the validated action's gain, the most favourable case for online learning that we found. A diagnostic showed that this learner repeats the action listed first, so its results depend on the order of the actions; an order-invariant version gains beyond the margin in none of $69$ cells. We report the diagnostic and the rule it adds to the protocol: permute the actions, or break ties at random.

\paragraph{Outline} Section~\ref{sec:related} places the work among parameter control, hyper-heuristics, algorithm selection and configuration, and benchmarking methodology. Section~\ref{sec:method} describes the hosts, instances, learners, the protocol and its estimators. Section~\ref{sec:results} answers the research questions and then tests how far the answers carry. Section~\ref{sec:discussion} discusses when online learning pays, gives a checklist for practitioners, and lists the limitations.
